\documentclass{article}
\usepackage[margin=1.15in]{geometry}
\usepackage{xcolor}
\usepackage{parskip}
\usepackage{fancyhdr}
\usepackage{array}
\usepackage{booktabs}
\usepackage{enumitem}
\usepackage{amsmath}
\usepackage{tabularx}
\usepackage{listings}

\pagestyle{fancy}
\fancyhf{}
\cfoot{\thepage}
\rhead{Lecture 13}
\lhead{MA 214: Applied Statistics (Spring 2026)}
\renewcommand{\headrulewidth}{.4mm}

\newcommand{\blankline}[1]{\underline{\hspace{#1}}}

\lstset{
  basicstyle=\ttfamily\small,
  columns=fullflexible,
  frame=single,
  framerule=0.4pt,
  breaklines=true,
  showstringspaces=false
}

\begin{document}

\noindent\textbf{Name:}\blankline{2.25in}\hfill\textbf{BUID:}\blankline{1.55in}

\medskip

\noindent\textbf{Name:}\blankline{2.25in}\hfill\textbf{BUID:}\blankline{1.55in}\newline

\begin{center}
 \Large In-Class Worksheet: Inference About the Slope
\end{center}

\subsection*{Scenario}
The Portal Project weighed $347$ kangaroo rats (\emph{Dipodomys merriami}) in 1999 --- the same animals as Lectures 9--10. We ask whether a rat's weight (g) rises with its hindfoot length (mm): \texttt{lm(weight \textasciitilde{} hindfoot\_length)}. The fit:

\begin{center}
\renewcommand{\arraystretch}{1.3}
\begin{tabular}{lrrrr}
\toprule
term & estimate & std.error & statistic & p.value \\
\midrule
\texttt{(Intercept)}     & $-12.68$ & $6.82$ & $-1.86$ & $0.064$ \\
\texttt{hindfoot\_length} & $1.585$  & $0.191$ & $8.31$ & $<0.001$ \\
\bottomrule
\end{tabular}
\end{center}

\newpage

\subsection*{Part 1: Set up and read the output}

\begin{enumerate}[label={\bfseries (\Alph*)},itemsep=.8em]
\item State $H_0$ and $H_A$ for the slope, in words. \\[3em]
\item Interpret the slope estimate $1.585$ in context, with units. \\[3em]
\item The intercept is $-12.68$ g. Is it meaningful here? Hindfoot lengths run from about $30$ to $58$ mm --- explain. \\[3em]
\item How many degrees of freedom does the $t$ statistic have? Give the number and the reason it is not $n-1$. \hfill $df = \blankline{1in}$ \\[1.5em]
\end{enumerate}

\subsection*{Part 2: The two roads for the test}

\begin{enumerate}[label={\bfseries (\Alph*)},itemsep=.8em]
\item \textbf{Road B (formula).} Confirm the reported $t$: compute $\hat\beta_1 / \mathrm{SE}(\hat\beta_1)$. \hfill $t = \blankline{1.2in}$ \\[1.5em]
\item A classmate shuffles the hindfoot--weight pairing $10{,}000$ times and finds that \emph{none} of the shuffled $|t|$ values reached the observed $8.31$. A simulated p-value is never $0$: use $(0+1)/(10{,}000+1)$, or report $p < 1/10{,}000$. How does that compare to the model-based $p < 0.001$? \\[3em]
\item In one sentence, why do the two roads agree here? What would have to be true of the data for them to \emph{disagree}? \\[3.5em]
\end{enumerate}

\newpage

\subsection*{Part 3: Find the bugs}
A classmate wants the shuffle p-value and the model p-value for the same test. \texttt{slope\_of(x, y)} correctly returns the slope. Their code runs without an error message --- and is wrong in \textbf{three} separate ways.

\begin{lstlisting}[language=R]
b_obs <- slope_of(kr$hindfoot_length, kr$weight)   # 1.585

# null distribution by shuffling
null <- replicate(10000,
  slope_of(sample(kr$hindfoot_length),
           sample(kr$weight)))

mean(null >= b_obs)              # randomization p-value

# model-based p-value, with SE = 0.191
t_obs <- b_obs / 0.191
2 * pt(-abs(t_obs), df = nrow(kr) - 1)
\end{lstlisting}

\begin{enumerate}[label={\bfseries (\Alph*)},itemsep=.8em]
\item Bug 1 is in the \texttt{null} line: they shuffled \emph{both} \texttt{hindfoot\_length} and \texttt{weight}. Why is shuffling one of them enough, and what should the line be? \\[3.5em]
\item Bug 2 is the randomization p-value line. The slope test is \emph{two-sided}, but \texttt{mean(null >= b\_obs)} counts only one tail. Write the correct line. \\[3.5em]
\item Bug 3 is the \texttt{df}. What should it be, and why does using $n-1$ matter less here (with $n = 347$) than it would for a small sample? \\[3.5em]
\end{enumerate}

\subsection*{Part 4: The interval, and the conditions}

\begin{enumerate}[label={\bfseries (\Alph*)},itemsep=.8em]
\item The model-based 95\% CI for the slope is $\hat\beta_1 \pm t^\ast\,\mathrm{SE}$ with $t^\ast \approx 1.97$. Compute it. \\[2.5em]
\item A bootstrap of the pairs gives a 95\% percentile CI of $[1.24, 2.06]$. Does it agree with (A)? It is not quite symmetric around $\hat\beta_1 = 1.585$ --- what feature of the data could make the bootstrap CI lean to one side? \\[3.5em]
\item You are handed the two diagnostic plots below in words: the residuals-vs-fitted plot shows a clear \emph{funnel} (spread widening to the right), and the Q--Q plot is roughly straight. Which LINE condition fails, which holds, and which road should you trust for this slope? \\[4em]
\end{enumerate}

\end{document}
